\addcontentsline{toc}{subsubsection}{Solution 9. The Reflected Product}

\begin{dlsummary}[title={\bfseries Solution 9. The Reflected Product},breakable]

\begin{proof}
\textbf{Step 1. Derive multiplicativity.}

The reciprocal identity at $x=1$ gives $f(1)=1$. Superadditivity shows that
$f$ is strictly increasing: if $v>u>0$, then
\[
f(v)\geq f(u)+f(v-u)>f(u),
\]
because $f(v-u)>0$. In particular,
\[
f(2)\geq f(1)+f(1)=2.
\]
Together with the given bound, this gives $f(2)=2$.

Apply the second inequality with $x$ replaced by $1/x$. Using the
reciprocal identity, we obtain
\[
\frac1{f(x)^2}\leq
\frac1{f(xy)f(x/y)},
\]
so $f(x)^2\geq f(xy)f(x/y)$. Hence
\[
f(x)^2=f(xy)f(x/y)\qquad(x,y>0). \tag{1}
\]
Taking $y=x$ in (1) gives $f(x^2)=f(x)^2$. Then, for any $a,b>0$,
use (1) with $x=\sqrt{ab}$ and $y=\sqrt{a/b}$ to get
\[
f(ab)=f(\sqrt{ab})^2=f(a)f(b).
\]
Thus $f$ is multiplicative.

\medskip
\noindent
\textbf{Method I. Dyadic squeeze.}

Multiplicativity and $f(2)=2$ give $f(2^n)=2^n$ for every integer
$n\geq0$. Fix $t>0$ and, for sufficiently large $n$, set
\[
m_n=\left\lfloor\frac{2^n}{t}\right\rfloor.
\]
Since $m_nt\leq2^n$, monotonicity and repeated superadditivity yield
\[
2^n=f(2^n)\geq f(m_nt)\geq m_nf(t).
\]
As $2^n/m_n\to t$, it follows that $f(t)\leq t$. The reciprocal
identity then gives
\[
f(t)=\frac1{f(1/t)}\geq t.
\]
Therefore $f(t)=t$ for all $t>0$.

\medskip
\noindent
\textbf{Method II. Dense fixed points.}

For every rational number $q=m/n$ with $m\in\mathbb Z$ and
$n\in\mathbb Z_{>0}$, multiplicativity implies
\[
f(2^q)^n=f(2^m)=2^m,
\]
where negative integer powers follow from $f(1)=1$ and
multiplicativity. Since $f(2^q)>0$, we get $f(2^q)=2^q$.
The set $\{2^q:q\in\mathbb Q\}$ is dense in $(0,\infty)$;
monotonicity and these fixed points imply $f(x)=x$ for every $x>0$.

Finally, $f(x)=x$ satisfies the reciprocal identity and $f(2)=2$;
both inequalities hold with equality. Hence it is the unique solution.
\end{proof}

\end{dlsummary}